\SourceLeaf{389}%
Aufnahme auf 1 zu 40000 und für den Stich auf 1 zu 80000. Danach veranschlagte man die Anzahl der Blätter (in derselben Größe wie oben) auf 208, die erforderliche Zeit auf 15 Jahre, die Kosten für die Arbeit, den Stich und den Druck von 3000 Exemplaren auf 4232000 Franken und schließlich den Verkaufspreis für jedes Blatt auf 7 Franken 50 Centimen oder für die ganze Karte auf 1560 Franken. Die gesamte Arbeit fällt in die Zuständigkeit des Depots. Man rechnet dabei aber mit der Mitwirkung der Katasterbehörde, obwohl aus dem Bericht nicht recht klar wird, in welchem Umfang. Wie es scheint, erwartet man von dieser (einem anderen Ministerium unterstellten) Behörde sämtliche Detailarbeiten, sodass den Ingenieurgeografen bei der Aufnahme nur die trigonometrischen Arbeiten und die Höhenbestimmungen zufallen. Diese Abhängigkeit von einer anderen Behörde, mit der es anscheinend keine recht harmonische Zusammenarbeit gibt (das Scheitern des ersten Plans wird der unzureichenden Unterstützung durch die Katasterbehörde zugeschrieben), könnte den erhofften raschen Fortgang dieses großen Unternehmens erheblich beeinträchtigen. Nach Abschluss der Arbeit soll noch ein umfangreiches Verzeichnis erscheinen, das sowohl die Zahlenwerte für Lage und Höhe der trigonometrischen Punkte als auch alle ihnen zugrunde liegenden Messungen bekannt macht. Damit würden freilich alle Wünsche erfüllt. Doch es ist zu befürchten, dass dieser Zeitpunkt noch sehr weit entfernt ist. Auch für weitergehende wissenschaftliche Fragen sind vor allem die Dreiecke und Dreieckspunkte erster Ordnung von größtem Interesse. Wir können daher den lebhaften Wunsch nicht unterdrücken, dass man die vollständige Veröffentlichung der Dreiecke erster Ordnung (die alle bereits gemessen sind) nicht so lange hinausschiebt und sie zum Nutzen der Wissenschaft sogleich zugänglich macht. Wer sich für höhere Geodäsie interessiert, wäre umso dankbarer, wenn die künftigen Bände des Memorials diese Wünsche erfüllten. da sie in den bisher erschienenen, in anderer Hinsicht so gehaltreichen Bänden am wenigsten berücksichtigt wurden und die wenigen theoretischen Abhandlungen und Hilfstafeln keinen Ersatz für die fehlenden \textit{Tatsachen} bieten, die das Depot doch so reichlich liefern könnte.
\begin{center}\rule{.3\textwidth}{.3pt}\end{center}
